\chapter{Basisuntersuchungen}

% \section{Untersuchungen können uns weiterhelfen}





% \section{Blutdruck- und Herzfrequenzmessung}

\label{topic:rr-messung}

% Grundlagen
% 
% Fallstricke
% 
% Die Blutdruckmanschette soll nach Möglichkeit nicht auf der Seite
% 
% \begin{itemize}
%     \item eines Shuntarmes bei Dialysepatienten
%     \begin{itemize}
%         \item Ein Dialyse-Shunt ist ein künstlicher Kurzschluss zwischen einer
%         Armarterie und einer Armvene. Dadurch wird die Vene erweitert und sie
%         pulsiert, mann kann viel mehr Durchfluss für die Dialyse erreichen.
%         Die Nahtstelle zwischen Vene und Arterie ist jedoch sehr sensibel, wenn
%         gestaut wird kann sie aufplatzen 
%         \item ${\rightarrow}$ Gefahr einer massiven Shuntblutung
%     \end{itemize}
%     \item eines Armödems
%     \item einer Brust-OP mit Lymphknoten-Entfernung (\"Odemgefahr)
%     \item einer Halbseitenschwäche oder -lähmung
% \end{itemize}
% 
% angelegt werden. 
% 
% \gCave{Shunt-Arme sind tabu für die Blutdruckmessung!}
% 
% Bitten Sie im Praktikum am KTW Dialysepatienten höflich Ihnen ihren
% Shunt-Arm zu zeigen!
% 
% Maßnahmen
% 
% \begin{enumerate}
%     \item Evt. Funktion des Blutdruckmessgerätes und des Stethoskops prüfen
%     \begin{itemize}
%         \item Kontrolle der Dichte des Ventils, der Verbindungsschläuche und
%         der Manschette (Ventil schließen, Manschette etwas aufpumpen und
%         Zusammenpressen, Manschette ganz entlüften)
%         \item Kontrolle der Gängigkeit des Ventils (es soll sich leicht auf
%         und zuschrauben lassen)
%         \item Kontrolle des Manometers auf korrekte Eichung: Die Nadel muss sich
%         im Ruhezustand genau bei 0\,mmHg befinden
%         \item Ohrstöpsel des Stethoskops zeigen schräg nach vorne
%         \item Trichterfunktion durch Berühren der Membran prüfen
%     \end{itemize}
%     \item Patientenlagerung
%     \begin{itemize}
%         \item Arm freimachen und entspannt auf einen Tisch legen in leichter
%         Flexion und Supination
%         \item Oberarm auf Herzhöhe
%     \end{itemize}
%     \item Auswahl der geeigneten Manschettengröße
%     \begin{itemize}
%         \item Manschette laut Herstellerangaben wählen z.B Standardmanschette:
%         Oberarmumfang max. 31 cm
%         \item Messwertverfälschung vermeiden! Der aufblasbare Teil der
%         Manschette soll 80\% des Oberarmumfangs umfassen, die Breite der
%         Manschette mindestens 1/3 des Oberarmumfangs betragen
%         \item Zu kleine Manschetten führen zum Messen zu hoher Werte, zu weite
%         Manschetten können nicht am Arm fixiert werden
%     \end{itemize}
%     \item Handhabung des Stethoskops
%     \item Bestimmen der Anlegestelle für Manschette und der
%     Auskultationsstelle durch Palpation
%     \begin{itemize}
%         \item Arteria brachialis medial in der Ellenbeuge aufsuchen
%     \end{itemize}
%     
%     \item Anlegen der Manschette
%     \begin{itemize}
%         \item Manschette muss auf unbekleideter Haut aufliegen
%         \item mäßig straff
%         \item Abstand zur Ellenbeuge 2,5 cm
%         \item aufblasbaren Teil über A. brachialis zentrieren
%         \item Schläuche nicht im Weg,
%         \item CAVE: Halbseitenschwäche/-lähmung, Shuntarm bei
%         Dialysepatienten, Armödem
%     \end{itemize}
%     \item Puls fühlen beim Aufpumpen
%     \begin{itemize}
%         \item Schnelles Aufpumpen
%         \item Tasten des Radialispulses mit Finger 2-4
%         \item Aufpumpen bis 30 mmHg über jenen Druck, ab dem kein Radialispuls
%         mehr getastet wird
%     \end{itemize}
%     \item Messvorgang
%     \begin{itemize}
%         \item Kopf des Stethoskops auf a. brachialis legen
%         \item Druck langsam ablassen 2 mmHg/sec
%         \item Systolischer Wert: Beginn der turbulenten Korotkoffschen
%         Strömungsgeräusche
%         \item diastolischer Wert: Verschwinden der Korotkoffschen Geräusche
%         \item Wenn Geräusche überhaupt nicht verschwinden (z.B. bei Kindern,
%         Schwangeren etc.), so gilt als diastolischer Wert der Messpunkt, an dem
%         die Geräusche deutlich leiser werden
%         \item Wenn Geräusche verschwunden sind, noch weitere 10mmHg langsam
%         ablassen -- wenn dann keine Geräusche mehr hörbar sind, kann die
%         Manschette zügig entleert werden
%         \item CAVE: auskultatorische Lücke
%     \end{itemize}
%     \item Protokollierung der Blutdruckwerte
% \end{enumerate}












\section{Sehen -- Hören -- Fühlen }
\gLayout{0}{Unsere eigenen Sinne}{
Die Seh-, Hör- und Tastsinne sind bei der körperlichen
Untersuchung unverzichtbar. Zuerst wird der Patient und sein
Körper betrachtet. Wir versuchen uns bereits bekannte Symptome wahrzunehmen. Hierbei können uns viele
augenscheinliche Sachen auffallen: zum Beispiel
\gE{Fettleibigkeit}, \gE{Zyanose}, \gE{Blässe},
\gE{Schweiß}, \gE{Fremdkörper}, \gE{Verletzungen}, \gE{Fehlstellungen},
\ldots

Der nächste Schritt betrifft das Hören: Hier sind zum
Beispiel die Atemgeräusche einzuordnen, die man eventuell schon ohne
Hilfsmittel vernehmen kann (brodelndes Atemgeräusch beim
Lungenödem, pfeifendes Atemgeräusch beim Asthmaanfall,
röcheln bei Atemwegsverlegung, \ldots)

Wenn wir gesehen und gehört haben können wir versuchen zu
fühlen: Wir können den (Radialis-) Puls fühlen, können
(stehende) Hautfalten beurteilen oder das Abdomen abtasten.
}{
\begin{itemize}
  \item Zyanose
  \item Blässe
  \item Schweiß
  \item Körperhaltung
  \item Fremdkörper
  \item Verletzungen, Fehlstellungen
  \item Atemgeräusche
  \item Puls
  \item Stehende Hautfalten
  \item Durchblutung des Nagelbetts
  \item \ldots
\end{itemize}
}{}{}{}




\section{Vitalwerte}
Das erheben von Vitalwerten ist eine Routinetätigkeit und
sollte bei jedem Notfallpatienten erfolgen. Es gibt nur wenig
Situationen bei denen man darauf verzichten kann (zum
Beispiel psychiatrische Patienten). Zu den Vitalwerten zählt man

\begin{itemize}
\item die \gEs{Herzfrequenz}, 
\item den \gEs{Blutdruck} und 
\item die \gEs{Atemfrequenz}.
\end{itemize}

Die Messung der Herzfrequenz erfolgt mittels Puls fühlen oder
mittels Pulsoxymeter. Beim Puls fühlen wird für
\gEs{15~Sekunden} der Puls gefühlt und die Pulsschläge
gezählt. Anschließend wird das Ergebnis \gEs{mit vier
multipliziert}, man erhält die Pulsschläge \gEs{pro Minute}
(\nicefrac{x}{min}). Neben der Frequenz wird auch die
Regelmäßigkeit beurteilt, das heißt ob der Puls \gE{rhythmisch} oder \gE{arrhythmisch} ist.

Alternativ kann die Messung mittels des \gE{Pulsoxymeters}
erfolgen (\prettyref{topic:pulsoxymetrie}). Dabei misst ein
Gerät die Herzfrequnz und zeigt diese an. Es ist jedoch zu
beachten, dass bei Notfallpatienten die Pulsoxymetrie oft
nicht funktioniert (z.B. Zentralisation bei Schock, kalte Finger,\ldots).

Die Messung des Blutdrucks ist unter \gSee{topic:rr-messung} beschrieben.

Bei der Messung der Atemfrequenz wird die Anzahl der Atemzüge pro Minute gezählt.\index{Atemfrequenz}

% TODO Atmung






\subsection{Messung der Herzfrequenz -- HF}
\label{topic:pulsmessung}

 
\gDefinition{Als \gT{Herzfrequenz} (\gT{HF})
bezeichnet man die Anzahl der Schläge des Herzens pro Minute.} 

Sie wird mittels Fühlen des Pulses
ermittelt. \marginpar{\gE{\scriptsize Die Blutgefäße werden im Detail unter
\gSee{topic:blutgefaesse-wichtige} besprochen.}} 
Dieser wird idealerweise an der
\gEs{Radialarterie} (A. \gE{radialis}) bzw. bei nicht
ansprechbaren oder schockierten Patienten an der
\gEs{Halsschlagader} (A. \gE{carotis}) ertastet. Bei
Neugeborenen und Säuglingen kann man auch an der
\gEs{Oberarmarterie} (A. \gE{brachialis}) fühlen. Gemessen
wird immer mit zwei bis drei Fingern. Nie jedoch mit dem Daumen, da hier meistens nur der eigene Puls gefühlt wird und nicht der des
Patienten. Wenn an der Halsschlagader (A. \gE{carotis}) die
Herzfrequenz ermittelt wird, darf nie auf beiden Seiten
gleichzeitig gemessen werden, da sonst die Gefahr besteht,
dass es zu einer Unterversorgung des Gehirn durch Blut kommen
kann.

\paragraph{Durchführung}

Es wird an der Radialarterie (A. \gE{radialis}) der Puls mit
zwei bis drei Fingern ertastet. Nun werden 15 Sekunden lang
die Schläge gezählt. Multipliziert man den Messwert mit vier, erhält man  die Herzfrequenz (HF) des Patienten. Der Normalwert liegt beim Erwachsenen zwischen 60 und 100 Schlägen pro Minute.
Dokumentiert wird sie mit einem HF (z.B.: HF\,=\,70)

\subsection{Messung des Blutdrucks -- RR}
\label{topic:blutdruckmessung}

Der Blutdruck ist jener Druck, welcher in unseren Blutgefäßen
herrscht. Wenn wir vom \emph{"`Blutdruck"'} sprechen beziehen wir
uns auf den Blutdruck in den
Arterien\footnote{\emph{Arterie:} Vom Herzen wegführendes
Blutgefäß.}. Er ist in Kombination mit der Herzfrequenz
einer der wichtigsten Messwerte im
Rettungsdienst.

Im Idealfall werden zwei Werte ermittelt: Der
\gEs{systolische Wert} ("`obere Wert"') ist jener Druck im
Blutgefäßsystem, der  während der \gE{Auswurfsphase} des Herzens, in der sich
die Herzkammern anspannen, zusammenziehen und Blut auswerfen, herrscht. Der
\gEs{diastolische Wert} ("`untere Wert"') ist jener Druck im
Blutgefäßsystem, der während der \gE{Entspannungsphase} des Herzens,
in der sich die Herzkammern mit
Blut füllen, herrscht. 

Die \gEs{Einheit}, in der der Blutdruck angegeben wird, ist
\gT{Millimeter Quecksilbersäule}
(\gT{mmHg}).\index{Millimeter Quecksilbersäule}\index{mmHg}\footnote{\gE{Hg} ist
das chemische Symbol für Quecksilber.}

\gMarried{
Der Blutdruck kann auf verschiedene Arten gemessen werden. 
%Üblich ist die
%\gE{indirekte Methode}, bei welcher mittels einer Druckmanschette an einer
%Extremität gemessen wird bei welchem Druck noch ein Blutfluß stattfindet. Dies
%erlaubt Rückschlüsse auf den Druck im Blutgefäß. 

Dazu gibt es unterschiedliche Messgeräte, grundsätzlich unterscheidet man
zwischen elektronischen und rein mechanischen Messgeräten.

\gEs{Elektronische Messgeräte}  werden oft von Patienten zur Selbstkontrolle des
Blutdrucks verwendet. Diese Geräte sind jedoch vergleichsweise fehleranfällig
und sind für einen professionellen Einsatz nicht geeignet. 

Für den
professionellen Einsatz gibt es auch spezielle
elektronische Blutdruckmessgeräte die zum Beispiel in Überwachungsmonitoren zum Einsatz
kommen, doch auch diese Geräte sind noch relativ fehleranfällig wenn die
Messbedingungen nicht optimal sind. 

% Der Blutdruck kann mittels elektronischer
% Blutdruckmessgeräte ermittelt werden. Diese sind aber im
% Rettungsdienst nicht in Verwendung, da es immer wieder zu
% Gerätefehler kommen kann. Außerdem können die Batterien leer
% werden. Solche Geräte haben meistens die Patienten für die
% tägliche Überwachung ihres Blutdrucks. Leider bedienen aber
% viele die Geräte falsch und fürchten sich dann
% lebensbedrohlich krank zu sein.

In der Praxis wird die Blutdruckmessung mit \gEs{mechanischen Messgeräten}
durchgeführt. Dabei wird eine aufblasbare Manschette mit angeschlossenem Dosenmanometer und
Pumpballen mit einem Ablassventil benötigt. Im Idealfall wird auch noch ein
Stethoskop verwendet. 
}{
\begin{center}
     
  \includegraphics[width=4cm]{./images/noauth/GER-rr-elektrisch-oberarm.png}
  
   \includegraphics[width=4cm]{./images/noauth/GER-rr-elektrisch-unterarm.jpg}
   
   \includegraphics[width=4cm]{./images/noauth/GER-rr-messgeraet.jpg}
\end{center}
 }

\subsubsection{Messung nach Riva-Rocci}\marginpar{\gZ{\scriptsize Benannt nach
dem Erfinder Scipione Riva-Rocci.}}

\gDefinition{Die Blutdruckmessung nach Riva-Rocci ist eine mechanische Messung mittels einer Druckmanschette.}

Umgangssprachlich wird die Blutdruckmessung nach Riva-Rocci als
\emph{"`RR-Messung"'}, bzw. der ermittelte Wert als \emph{"`RR"'} bezeichnet.

Bei der RR-Messung gibt es zwei Möglichkeiten zu
einem Wert zu kommen. Die genauere ist die
\gEs{auskultatorische Methode} (durch hören von
Strömungsgeräuschen). Die zweite Möglichkeit ist die
\gEs{palpatorische Methode} (durch tasten des Pulses). 

\gEs{Der diastolische Wert kann nur
bei der auskultatorischen Methode ermittelt werden!} 

\paragraph{Durchführung}

\subparagraph{\gT{Auskultatorische Methode}}

Die Manschette sollte faltenfrei und nicht über Kleidung angelegt werden. Die
Mitte der Luftkammer soll über der Oberarmarterie zu liegen kommen. Manche
Manschetten haben einen Markierungspfeil für die Oberarmarterie, der Pfeil soll
dann oberhalb dieser zu liegen kommen. Man achte auf die Lage der Innen- bzw.
Außenseite!

Nun wird
der Puls an der Radialarterie (A. radialis) getastet. Die Manschette wird
solange aufgepumpt, bis der Puls nicht mehr getastet werden kann. Dann werden
noch 20-30\,mmHg weiter aufgepumpt. Das Stethoskop wird in der Ellenbeuge auf
die Oberarmarterie (A. brachialis, Verlauf eher zur Körpermitte) angelegt. Die
Luft wird langsam aus der Manschette abgelassen. 

Der Wert bei dem man
erstmalig ein Pochen \gZ{(Korotkowsche Geräusche)} hört gilt als
\gEs{systolischer Wert}. Die Luft wird weiter abgelassen. Der Wert, bei dem das
Pochen verschwindet, ist der \gEs{diastolische Wert}. \footnote{Es kann
vorkommen dass das Pochen nie verschwindet sondern nur plötzlich leiser wird.
Dann gilt der Wert bei dem das Pochen \emph{deutlich} leiser wird als
diastolischer Wert.}

Die
Werte werden auf die nächste Fünferstelle abgerundet (z.B.:
122\,/\,82\,mmHg wird abgerundet auf 120\,/\,80\,mmHg). Dokumentiert
wird der Blutdruck mit RR (z.B.: RR\,=\,125\,/\,75\,mmHg). Der
Normalwert beim Erwachsenen liegt ca. bei RR 120\,/\,80\,mmHg
(\gSee{topic:blutdruck}).


\begin{figure}[H]
    \begin{center}
   
   
    \includegraphics[width=10cm]{./images/noauth/GER-rr-geraeusche.png}

   
    \gCaptionof{figure}{\gAuthNotesPicLegalNotOk{GER-rr-geraeusche.png}{}Prinzip
    der auskultatorischen RR-Messung: In Abhängigkeit vom außen anliegenden
    Manschettendruck kommt es zu einem unterschiedlich starken Blutfluss und dadurch entstehenden, hörbaren, Strömungsgeräuschen.}
    \end{center}
\end{figure}




\subparagraph{\gT{Palpatorische Methode}}
Die Manschette sollte faltenfrei, nicht über Kleidung, in der
Mitte des Oberarms angelegt werden. Nun wird der Puls an der
Radialarterie (A. radialis) ertastet. Die Manschette
wird solange aufgepumpt bis der Puls nicht mehr getastet werden
kann. Nun noch 20\,mmHg weiter aufpumpen. Die Finger bleiben
auf der Radialarterie (A. radialis). Die Luft langsam
aus der Manschette ablassen. Sobald man wieder einen Puls fühlt,
muss man den Messwert am Dosenmanometer ablesen. Dies ist der
\gEs{systolische} Wert. 

\gE{Der diastolische Wert kann mit
dieser Methode nicht ermittelt werden.} In der
Notfallmedizin ist der systolische Wert jedoch ohnehin in der
Regel der wichtigere.

\paragraph{Was ist besonders zu beachten?}

\marginpar{\includegraphics[width=\linewidth]{./images/sign_cardive.jpg}}

Die Blutdruckmanschette soll nach Möglichkeit nicht auf der Seite
\begin{itemize}
	\item eines Shuntarmes bei Dialysepatienten
	\begin{itemize}
		\item Ein Dialyse-Shunt ist ein künstlicher Kurzschluss
		zwischen einer Armarterie und einer Armvene. Dadurch wird
		die Vene erweitert und sie pulsiert. Mann kann viel mehr
		Durchfluss für die Dialyse erreichen. Die Nahtstelle
		zwischen Vene und Arterie ist jedoch sehr sensibel, wenn
		gestaut wird kann sie aufplatzen $\rightarrow$ Gefahr
		einer massiven Shuntblutung.
	\end{itemize}
		\item eines Armödems, z.B. nach einer
		Brustamputation
		\item einer Halbseitenschwäche oder -lähmung 
\end{itemize}
angelegt werden.

\gCave{Shunt-Arme sind tabu für die Blutdruckmessung!}









\section{Apparative Untersuchungen}

\subsection{Pulsoxymetrie}
\label{topic:pulsoxymetrie}



\gMarriedNG{\gTxBeschreibung}
{
Die Pulsoxymetrie ist eine einfache apparative
Untersuchungsmethode, welche mithilfe eines
\gE{Pulsoxymeters} durchgeführt wird. 
Sie dient zur nicht-invasiven Überwachung der
\gT{Sauerstoffsättigung} (\gT{Sp\ce{O2}}) des kapillären
Blutes. Dies geschieht mittels Photorezeptoren und
Infrarotlicht, mit deren Hilfe das Verhältnis der mit
Sauerstoff beladenen roten Blutkörperchen zu den nicht
beladenen bestimmt wird. 

Als weiteren Wert zeigt das Pulsoxymeter die
\gEs{Pulsfrequenz} an.
\cite{LoefflerBiochemie7,ArbeitstechnikenAZ1}

}{
% TODO Slide fehlt.
{\gSymbolText}
}

\gLayout{2}{Durchführung}
{
Es wird ein eine Klemme mit einer besonderen Lichtquelle
und einem Sensor auf den Finger des Patienten
geklipt. Das Licht wird vom
gegenüberleigenden Sensor aufgefangen, das Gerät kann aus den daraus gewonnenen Informationen auf den Sauerstoffgehalt des Blutes und die Pulsfrequenz schließen. 

Dies setzt jedoch
voraus das der Patient in der Peripherie über eine
ausreichende Durchblutung verfügt, da das Gerät sonst keine
Informationen gewinnen kann (mangels Blut zur
Beurteilung). 
}{}{}{}{}


\gFazit{Eine ausreichende Blutversorgung der Körperstelle,
an der der Sensor angebracht ist, ist Voraussetzung für die
Pulsoxymetrie.}

Das Pulsoxymeter eignet sich sehr gut für das Monitoring,
ersetzt aber niemals das wachsame Auge auf den Patienten!

\gLayout{0}{Ermittelte Werte}
{
Wie Sauerstoffsättigung wird in Prozent
angegeben und wird mit \gT{Sp\ce{O2}} abgekürzt. Der Normalwert
liegt zwischen \gEs{96-100\,\%}. Bei Patienten mit COPD
findet man jedoch oft sehr niedrige Werte (z.B.
90\,\%), ohne dass die Patienten Beschwerden haben. }{
% TODO Slide fehlt.
% \begin{itemize}
%     \item \gZ{Point}
% \end{itemize}
\begin{itemize}
  \item Normalwert: 96-100\, \%
  \item COPD-Patient: auch Werte < 90\, \% möglich ohne
  Beschwerden
\end{itemize}
}{}{}{}


\gMarriedNG{Beurteilung}
{
Bei der Interpretation des
$SpO_2$ muss man sehr vorsichtig sein! Auch ein Normalwert bedeutet nicht, dass der Patient
genug Sauerstoff hat beziehungsweise, dass keine
Sauerstoffgabe erfolgen soll. Die Sauerstoffgabe richtet
sich immer nach dem klinischen
Zustandsbild des Patienten und dessen Diagnose, und nur in
Ausnahmefällen (z.B. Asthma bronchiale,
\prettyref{topic:asthma-bronchiale}) nach dem
Sp\ce{O2}-Wert. Der Sp\ce{O2}-Wert soll immer im Verlauf beurteilt werden. }{
% TODO Slide fehlt.
% \begin{itemize}
%     \item \gZ{Point}
% \end{itemize}
{\gSymbolText}
}

\gFazit{Behandle den Patienten und nicht das Gerät!}

Immer
zuerst auf den Patienten schauen (Symptome) und erst dann auf das Pulsoxymeter!



\gLayout{0}{\gTxGefahren{} und Fehlerquellen}
{
Eine besondere Fehlerquelle ist die
Kohlenmonoxid-(\ce{CO})-Vergiftung, bei ihr wird nicht nur
der Patient schweinchenrosa obwohl er eigentlich zu wenig
Sauerstoff hat, auch das Pulsoxymeter verwechselt
Kohlenmonoxid mit Sauerstoff und zeigt falsch hohe Werte an!

Voraussetzung für eine korrekte Messung ist eine ausreichende
Durchblutung der entsprechenden Meßstelle. Z.B. durch Kälte,
Schock oder auch eine Blutdruckmessung kann die Messung
gestört oder unmöglich gemacht werden. Weitere Hindernisse
sind z.B. kalte Finger (verringerte Durchblutung) oder
Nagellack (Licht kann den Nagellack nicht durchdringen,
bzw. dessen Farbe wird verfälscht.) Mit
Nagellack-Entferner kann der Nagel eventuell
gereinigt werden. Alternativ kann der Sensor auch
um 90\textdegree gedreht werden. 

Von der Seite einstrahlendes Licht kann ebenso das Ergebnis
verfälschen, da es sich ja um eine lichtabhängige
Meßmethode handelt.
}{
\begin{itemize}
	\item \ce{CO}-Vergiftungen (falsch hohe Wert)
	\item Minderdurchblutung der Peripherie; Entweder durch Zentralisation beim Schock oder Unterkühlung.
	\item Nagellack
	\item Künstliche Fingernägel
	\item Sonneneinstrahlung von der Seite
\end{itemize}
}{}{}{}




\subsection{Blutzuckermessung}
\label{topic:blutzucker-messung}

\gLayout{0}{Interpretation und Normalwert}
{
Der Blutzuckerspiegel wird in Milligramm pro Deziliter
(\nicefrac{mg}{dl}\gZ{, enspricht mg\,\%}) gemessen. Der
Wert ist abhängig von der letzten Mahlzeit des Patienten.
Nach der Nahrungsaufnahme steigt der Blutzuckerspiegel
stark an und fällt dann wieder auf einen "`nüchternen
Normalwert"' ab. Der normale
\gEs{Nüchtern-Blutzuckerspiegel} bewegt sich zwischen
\gEs{80\,--\,120\,\nicefrac{mg}{dl}}. Im Rettungsdienst
kann der Normalbereich jedoch großzügiger angesetzt werden,
da bei Werten zwischen 60 und 200 \nicefrac{mg}{dl} keine
Notfallsymptome auftreten.

}{
\begin{itemize}
    \item Abhängig von letzter Mahlzeit!
    \item Normalwert: 80\,--\,120\,\nicefrac{mg}{dl}
    (nüchtern)
    \item Bereich ohne Notfallsymptome:
    60\,--\,200\,\nicefrac{mg}{dl}
%     \item Notfallrelevanter Hypoglykämie: <
%     60\,\nicefrac{mg}{dl}
%     \item Notfallrelevante Hyperglykämie: >
%     200\,\nicefrac{mg}{dl}
\end{itemize}
}{}{}{}

\gMarriedNG{Material \& Personal}
{
\begin{enumerate}
  \item 1 eingeschulte Person
  \item Kontamed-Box
  \item Alkoholtupfer
  \item Trockene Tupfer
  \item Blutzuckermessgerät
  \item Blutzuckermessstreifen
  \item Lanzette
\end{enumerate}
}
{}


\gMarriedNG{Durchführung}
{
%\begin{multicols}{2}
\begin{enumerate}
  \item Finger auswählen: möglichst von der nichtdominanten
  Hand.
  \item Einstichstelle bestimmen: Am Fingerendglied,
  seitlich (Rand!), nicht in die Fingerkuppe, nicht zu nahe
  am Gelenk. Außenseite des Ringfinger bevorzugen.
  \item Stelle desinfizieren: Mit Alkoholtupfer abwischen
  und gemäß Einwirkzeit einwirken lassen.
  \item Blutzuckermessstreifen halb in das Gerät schieben,
  aber noch nicht einrasten lassen.
  \item Mit einer schnellen, beherzten Bewegung den
  Patientenfinger mit der Lanzette stechen und sofort
  herausziehen \gE{("`Rein-raus-Methode"')}.
  \item Lanzette in die \gE{Kontamed-Box} entsorgen.
  \item Ersten Bluttropfen mit trockenem Tupfer abwischen.
  \item Messstreifen ganz in das Gerät schieben, das Gerät
  schaltet sich ein und führt einen Selbsttest durch. Wenn
  im Display ein Tropfen angezeigt wird ist das Gerät
  messbereit.
  \item Den Messstreifen mit dem Gerät in den Bluttropfen
  tauchen, das Blut steigt auf\footnote{Kapillarwirkung}.
  \item Beim Piepton ist genug Blut in die Testkammer
  eingeströmt und das Gerät beginnt mit der Messung und kann
  entfernt werden. 
  \item Mit einem trockenen Tupfer auf die Einstichstelle
  drücken (lassen).
  \item Ein neuerlicher Piepton zeigt an, dass die Messung
  beendet wurde und der Wert vom Display abgelesen
  werden kann. 
  \item Messstreifen entsorgen\footnote{Entsorgung gemäß
  Abfallentsorgungsplan. \opt{ORGASBFLO}{\ASBFLO{} Der
  Meßstreifen kann in den normalen Müll entsorgt werden.}}.
\end{enumerate}
%\end{multicols}
}{

}




\subsection{Körpertemperatur}

Die Körpertemperatur kann mit unterschiedlichen Mitteln und
an unterschiedlichen Stellen bestimmt werden. 

\gMarriedNG{Messarten}
{
Mit dem
\gEs{herkömmlichen Fieberthermometer} kann die Temperatur
unter der \gE{Achsel} (\gE{axillär}) oder im \gE{Enddarm}
(\gE{rektal}, v.a. bei Kleinkindern) bestimmt werden. Die
Messung im Mund ist in der Praxis eher unüblich. Es ist bei der Messung eine
Schutzhülle zu verwenden. 

Mit dem \gEs{Ohrthermometer} kann mittels Infrarot die
Temperatur im \gE{Mittelohr} bestimmt werden. Es sind die
entsprechenden \gE{Schutzkappen} zu verwenden.

Die Temperatur im Enddarm und im Mittelohr bezeichnet man
als \gT{Körperkerntemperatur}, diese beträgt um die
37\textcelsius. Die unter der Achsel gemessene Temperatur
kann ca. \nicefrac{1}{2}\textcelsius{} niedriger sein.
}{
% TODO Slide fehlt.
% \begin{itemize}
%     \item \gZ{Point}
% \end{itemize}
{\gSymbolText}
}


\section{Körperliche Untersuchungen}

\subsection{Abtasten des Abdomens}

\gLayout{2}{Abtasten des Abdomens}
{\label{topic:untersuchung-abdomen-abtasten}
Das Abdomen kann man abtasten um die Anspannung der
Bauchdecke zu beurteilen (sie voraus), Schmerzen zu
beurteilen und zu zuordnen, Fremdkörper zu ertasten und
anderes. Man beginnt damit den Patienten zu fragen ob er ihm
auch Schmerzen hat und wo er diese hat. Die Untersuchung
beginnt man in dem Quadranten der am weitesten von den
Schmerzen entfernt ist!

Man legt nun eine Hand flach auf den Bauch des Patienten und
die andere Hand flach auf die erste. Mit den Fingern drückt
man nun vorsichtig mehrere Zentimeter auf die Bauchdecke und
beurteilt, ob der Bauch weich oder verhärtet ist. Dabei
fragt man den Patienten ob er Schmerzen empfindet. Dies
wiederholt man in allen vier Quadranten.

Weiters kann man beurteilen ob der Patient Schmerzen durch
den Druck hat (Druck Schmerz) oder ob er Schmerzen beim
loslassen bekommt
(\gT{Loslassschmerz}\index{Loslassschmerz}).
}
{
% TODO VIT: Slide fehlt!
% \begin{itemize}
%     \item \gZ{Point}
% \end{itemize}
{\gSymbolText}
}{}{}{}

\subsection{Neurocheck}

%\begin{multicols}{2}

\gLayout{1}{}{
\begin{enumerate}
  \item \gEs{Bewusstseinslage}: Für eine grobe Einschätzung
  ist folgende Differenzierung ausreichend:
  	\begin{enumerate}
        \item bewusstseinsklar
        \item bewusstseinsgetrübt
        \item bewusstlos
      \end{enumerate}
      Notärzte und Notfallsanitäter verwenden zur
      Beurteilung des Bewusstseins eine detailliertere
      Skala, die so genannte Glasgow Coma Scale.
  \item \gEs{Orientierung}:
  	\begin{enumerate}
        \item zur Person (Patient kennt seinen Namen)
        \item zum Ort (Patient weiß wo er sich befindet)
        \item zur Zeit (Patient kann das ungefähre Datum
        und die ungefähre Uhrzeit nennen)
    \end{enumerate}
  \item \gEs{Pupillen} Begutachtung und Lichtreaktion
\begin{enumerate}
  \item Pupillen gleich weit?
  \item Blickstarre (\gT{Herdblick})? \index{Herdblick}
  \item \gEs{Lichtreaktion:} In jedes Auge 2\,x leuchten und
  jeweils das gleiche und das andere Auge beobachten
  
  Beurteilung von:
  \begin{enumerate}
      \item Orientierung zu Person, Ort, Zeit und Situation.
  \item Schnelligkeit der Reaktion: prompt, verzögert?
  \item Reagieren immer beide
  Pupillen?~\gZ{\footnote{\gZ{isokor}}}
    \end{enumerate}
\end{enumerate}
  \item Lichtscheu und Sehstörungen?
  \item \gEs{Kraftprobe} grobe Prüfung an oberer und
  unterer Extremität, beurteilt wird die Kraft und die
  Seitengleichheit. 
    \begin{enumerate}
  \item Patient überkreuzt die Hände geben, bei kräftigen
  Patienten nur Zeige- und Mittelfinger.
  \item Patienten bitten kräftig zuzudrücken.
  \item Mit Handflächen von unten gegen die Zehen drücken
  und Patienten bitten \emph{"`auf das Gaspedal zu
  drücken"'}.
  \item Von oben auf die Zehen drücken und Patienten bitten
  die Zehen Richtung Kopf zu heben.
\end{enumerate}
  \item \gEs{Nackensteifigkeit} (\gT{Meningismus}) Der
  Patient liegt flach am Rücken. Der Untersucher fasst den
  Kopf des Patienten und versucht das Kinn an die Brust zu
  führen. Der Patient darf dabei nicht mithelfen.
  Normalerweise ist diese Bewegung ohne Probleme möglich.
  
  Wenn eine Nackensteifigkeit vorliegt ist die Beugung nicht
  möglich und der Versuch schmerzhaft, der Patient bewegt
  evt. den ganzen Rumpf mit. Die Bedeutung der
  Nackensteifigkeit wird unter \gSee{topic:meningitis}
  besprochen. Bei jedem Patienten mit Fieber ist auf das
  Vorliegen einer Nackensteifigkeit zu untersuchen.  
  \item \gEs{Blutzuckermessung} Eine Blutzuckermessung soll
  grundsätzlich bei einem Neurocheck erfolgen.
\end{enumerate}
}{
\begin{itemize}
  \item Bewusstseinslage
  \item Orientierung
  \item Pupillen
  \item Lichtscheue, Sehstörungen
  \item Kraftprobe
  \item Nackensteifigkeit
  \item Blutzuckermessung
\end{itemize}
}{}{}{}

%\end{multicols}

\gLayout{57}{}
{}{}
{
\begin{tabularx}{\linewidth}{XXX}
\includegraphics[width=\linewidth]{./images/motal-aass/neuro1.jpg}
&
\includegraphics[width=\linewidth]{./images/motal-aass/neuro2.jpg}
&
\includegraphics[width=\linewidth]{./images/motal-aass/neuro3b.jpg}
\\
Pupillenreaktion, Lichtscheue, Sehstörungen 
&
Kraftprobe an den Armen
& 
Kraftprobe an den Beinen
\\\addlinespace
\includegraphics[width=\linewidth]{./images/motal-aass/neuro4.jpg}
&
\includegraphics[width=\linewidth]{./images/motal-aass/neuro5.jpg}
&
\includegraphics[width=\linewidth]{./images/motal-aass/neuro6.jpg}
\\
Nackensteifigkeit
&
Nackensteifigkeit
&
Blutzuckermessung
\\\addlinespace\bottomrule
\end{tabularx}
}{Bilderserie: Neurocheck}{}

% FEHLT: Bild für Orientierung (stattdessen jetzt 2x
% Nackensteifigkeit)

\subsection{Traumacheck}
\label{topic:traumacheck}
\gDefinition{Suche nach Verletzungen oder Verletzungsanzeichen am (liegenden)
Patienten.}

\gLayout{0}{Ausführlicher Traumacheck}
{
Der
hier beschriebene Traumacheck ist ein "`ausführlicher"' (vollständiger)
Traumacheck. In der Praxis werden beim Patienten zur raschen Einschätzung manche Punkte zuerst
gemacht und andere Schritte auf später verschoben (\gE{schneller} vs.
\gE{ausführlicher} Traumacheck).

\begin{itemize}
\item Patient orientiert? Unfallhergang erinnerlich? Schmerzen?

\end{itemize}
\begin{itemize}
    \item Abtasten mit festem Griff, \gEs{bei Schmerzen die vom Patienten erwähnte
    Stelle als letztes untersuchen}: von Kopf bis Fuß
    \begin{itemize}
        \item Schädel: Schädeldach, Hinterkopf,
        Stirn, Wangenknochen, Nasenbein, Kiefer.
        \item Hals.
        \item Obere Extremität: Schultern, Arme, Hände.
        \item Brustkorb: vorne und seitlich (Prellmarken? Atemmechanik?, gegen leichten
        Druck einatmen lassen).
        \item Bauch in 4 Quadranten (Prellmarken?
        Abwehrspannung/harte Bauchdecke?),
        \item Becken (stabil?), Schambein.
        \item Untere Extremität: Beine (sichtbare Verletzungen?), Füße:
        Schuhe und Socken ausziehen)
    \end{itemize}
    \item Durchblutung, Motorik, Sensibilität (\gT{DMS})
    \begin{itemize}
        \item Finger und Zehen: Kontrolle von \gE{Motorik} (bewegen lassen), \gE{Sensibilität} ({\quotedblbase}Spüren
        Sie das?{\textquotedblleft}) und \gE{Durchblutung}:
        
        \item \gT{Nagelbettprobe} (\gEs{Rekapillarisierungszeit}): leichter Druck auf
        Nagelende bis Nagelbett weiß ist, loslassen. Zeit bis Nagelbett wieder
        rosarot ist sollte nicht länger als 1-2 Sekunden sein.
        Seitenvergleich! Eine verlängerte Rekapzeit deutet auf eine
        Durchblutungsstörung hin. (einseitig oder beidseitig? lokal oder systemisch?)
    \end{itemize}
    \item Kontrolle von Mund, Nase und Ohren auf Blutungen (mit einer
    Lampe).
    
    \item \gT{Tupfertest}\index{Tupfertest}\label{topic:tupfertest}: etwas Blut aus
    Ohren oder Nase mit einem Tupfer auffangen. Sollte Liquor im Blut enthalten sein bildet sich um den roten Blutfleck
    ein heller, bernsteinfarbener Hof.
    
\end{itemize}
Bei Hinweis auf eine Schädelverletzung ist ein Neurocheck auf jeden Fall
erforderlich!
}{
\begin{itemize}
  \item Bewusstsein, Orientierung
  \item Abtasten
  \begin{itemize}
    \item Schmerzen
    \item Blutungen
    \item Fehlstellungen
    \item DMS
    \item Atemmechanik
    \item harte Bauchdecke
  \end{itemize} 
  \item Kontrolle von Mund, Nase, Ohren
  \item Tupfertest
\end{itemize}
}{}{}{}

\gLayout{57}{}
{}{}
{
\begin{tabularx}{\linewidth}{XXX}
\includegraphics[width=\linewidth]{./images/motal-aass/traumacheck1.jpg}
&
\includegraphics[width=\linewidth]{./images/motal-aass/traumacheck2.jpg}
&
\includegraphics[width=\linewidth]{./images/motal-aass/traumacheck3b.jpg}
\\
Kleidung entfernen 
&
Schuhe ausziehen (Fixierung, Stiefelgriff) 
& 
Untersuchung des Kopfs 
\\\addlinespace
\includegraphics[width=\linewidth]{./images/motal-aass/traumacheck4.jpg}
&
\includegraphics[width=\linewidth]{./images/motal-aass/traumacheck5b.jpg}
&
\includegraphics[width=\linewidth]{./images/motal-aass/traumacheck6b.jpg}
\\
In alle Körperöffnungen im Kopf leuchten, Pupillenkontrolle.
&
Neben Blut kann auch Liquor austreten (Tupfertest).
&
Untersuchung des Thorax (inkl. Atemmechanik)
\\\addlinespace
\includegraphics[width=\linewidth]{./images/motal-aass/traumacheck7b.jpg}
&
\includegraphics[width=\linewidth]{./images/motal-aass/traumacheck8b.jpg}
&
\includegraphics[width=\linewidth]{./images/motal-aass/traumacheck9b.jpg}
\\
Untersuchung des Abdomen und des Beckens
& 
Untersuchung der Extremitäten (DMS, Nagelbettprobe) 
&
Nach dem Traumacheck den entkleideten Patient vor
Wärmeverlusten schützen. 
\\\addlinespace\bottomrule
\end{tabularx}
}{Bilderserie: Traumacheck}{}

% }
% {
% {\gSymbolText}
% }

\gMarriedNG{Schneller Traumacheck}
{
Um rasch zu erkennen ob der Patient vital bedroht ist werden in der Praxis
wichtige Schritte vorgezogen.
\begin{itemize}
    \item Abtasten von

\begin{enumerate}
    \item Kopf,
    \item Hals,
    \item Thorax inkl. Atemprobe,
    \item Bauch,
    \item Becken und 
    \item Oberschenkel.
\end{enumerate} 
\item Kontrolle von Mund, Nase und Ohren auf Blutungen (mit einer
    Lampe).
\item Kann der Patient \gE{Hände} und \gE{Füße} spüren und bewegen?
% KOCH: komplettes DMS würd zu lange dauern, Schuhe ausziehen, etc. 
% GABS: Rekap-Zeit ist eh im normalen Ersteinschätzungsblock dabei.
% MOTM: POI 35: "`ja klingt verünftig."'
\end{itemize}
}
{
\begin{itemize}
    \item Rasche Ersteinschätzung.
    \item Rasches Erkennen von schweren, lebensbedrohlichen Verletzungen.
\end{itemize}
{\gSymbolText}
}




